\chapter{Finite Fusion with Exact Statistical Inheritance}
\label{ch:fusion}

This chapter constructs a population whose dynamics makes agents coincide at finite times and whose coincident blocks carry exact inherited laws, evaluator, metric, and responses. The dynamics is a nonsmooth modification of the reciprocal law potential of \cref{def:dyn-potential}, made by four constitutive changes stated in \cref{sec:fus-dynamics}: each peer divergence is replaced by the square root of twice a divergence, the receiver's law is restored to unit log variance before it is compared, the peer term carries a multiplicity-weighted coefficient, and source rows are fixed. The laws are restricted to a lognormal family with fixed variances. The potential is not, on any open set of means that meets a coincidence hyperplane $\{\mu_i^\alpha=\mu_j^\alpha\}$, the free energy of a normalized model whose log density is integrable against these laws: in logarithmic coordinates such a free energy is, up to a constant, the convolution of the negative log density with a fixed Gaussian kernel, hence real analytic in the means, whereas the potential has a corner on every coincidence hyperplane. Its conclusions do not transfer to the smooth potential or to the per-agent and global free energies of \cref{ch:dynamics}.

Pairwise fusion penalties occur in convex clustering \cite{Hocking2011}, and suitable nonsmooth consensus protocols produce finite-time agreement \cite{Cortes2006}. On the restricted family the square-root interaction is exactly the absolute-value kernel. The population is a weighted particle solution of the two-dimensional aggregation equation with kernel $\lambda\|x\|_1$ in a quadratic external potential; on the common-offset family of \cref{fus:sticking}, after a rescaling of positions and a change of time, each sector is the one-dimensional aggregation equation with kernel $\lambda|z|$, whose sticky collisions and finite-time total collapse are known \cite{CarrilloEtAl2011,BonaschiEtAl2015}. What the construction adds is the inheritance bookkeeping: represented laws, model evaluator, transport, parent metric, and boundary kernels descend exactly through each collision. Its dynamical outcome is a finite coalescence tree that ends in global consensus for every initial datum (\cref{fus:sync}); persistent heterogeneous organization does not occur. On the common-offset family the tree is a linkage dendrogram of the initial data that does not depend on the coupling strength (\cref{fus:ncl}). \Cref{sec:fus-status} reports the construction against the formation standard of \cref{inh:recursive-formation}.

\section{A law family closed under geometric averaging}

Work over one context, with state and model spaces
$\mathsf K=\mathsf M=(0,\infty)$. Write
$\operatorname{LN}(\mu,\sigma^2)$ for the law of $e^Z$ when
$Z\sim N(\mu,\sigma^2)$. An agent of integer multiplicity $n\geq1$ carries
the product law
\begin{equation}
 q_n=\operatorname{LN}(\mu^q,1/n),\qquad
 s_n=\operatorname{LN}(\mu^s,1/n),\qquad
 \zeta_n=q_n\otimes s_n.
 \label{fus:law-family}
\end{equation}
Multiplicity is retained state information. The variances are fixed within
each type; only the two log means evolve. This is a constraint on the
dynamics, not an invariant family of a larger flow. If the common log
variance $v$ were released, the local divergence of
\cref{fus:local-inheritance} below would have derivative $n(a-1)$ in $v$ at
$v=1/n$, and the peer comparison between restored messages would have a
nonzero source-variance derivative whenever the means differ. The Fisher
flow with free variances therefore leaves the family, and releasing the
variance destroys the encoding of multiplicity in the law. Laws are first
written in a common reference coordinate. A fixed positive frame
$U_i^q=U_i^s=:U_i\in\GL^+(1)$ represents both sectors locally by
multiplication. Frames do not evolve in this example.

Choose constants $a>c>0$ and define the normalized reference and evaluator
\begin{equation}
 r_n=\operatorname{LN}\left(0,\frac{a}{n(a^2-c^2)}\right),\qquad
 G_n(m,dk)=\operatorname{LN}\left(\frac ca\log m,\frac1{an}\right)(dk).
 \label{fus:evaluator}
\end{equation}
The two occurrences of a lognormal expression in this equation denote
ordinary probability laws and kernels. In logarithmic coordinates the joint
reference $r_n(dm)G_n(m,dk)$ has precision
\begin{equation}
 nA,\qquad A=\begin{pmatrix}a&-c\\-c&a\end{pmatrix}.
 \label{fus:precision}
\end{equation}
Its eigenvalues $n(a-c)$ and $n(a+c)$ are positive. The predictive law is
\begin{equation}
 s_nG_n=\operatorname{LN}\left(
 \frac ca\mu^s,\frac{(c/a)^2+1/a}{n}\right).
 \label{fus:predictive-law}
\end{equation}
Thus the evaluator depends on the represented model value. Physical
references and evaluators are obtained by pushing both coordinates through
the common local frame.

\propositionheading{Local energy and metric}{fus:local-inheritance}
For every $n\geq1$ and every pair of real means, let
\begin{equation}
 C=a-1-\frac12\log(a^2-c^2),\qquad
 \mathcal U(\mu^q,\mu^s)
 =\frac a2\bigl((\mu^q)^2+(\mu^s)^2\bigr)-c\mu^q\mu^s.
\end{equation}
Then
\begin{equation}
 E_{0,n}=\KL(q_n\otimes s_n\|r_nG_n)
 =C+n\mathcal U,\qquad
 g_n=n\bigl((d\mu^q)^2+(d\mu^s)^2\bigr).
 \label{fus:local-formulas}
\end{equation}
\begin{proof}
The logarithm is a measurable bijection, so it preserves KL divergence.
The recognition law in log coordinates has covariance $I_2/n$, while the
reference has covariance $A^{-1}/n$. The Gaussian KL formula gives
\begin{equation}
 E_{0,n}=\frac12\left[
 \operatorname{tr}A-2-\log\det A
 +n(\mu^q,\mu^s)A(\mu^q,\mu^s)^{\mathsf T}\right],
\end{equation}
which is the stated expression. The score for either mean is
$n(Z-\mu)$, whose squared expectation is $n$. The product law has zero
cross score expectation, giving the displayed Fisher metric.
\end{proof}

For independent synchronized children of multiplicities $n_j$, set
$n_I=\sum_{j\in I}n_j$. Their retained reference samples are
\begin{equation}
 k_I=\prod_{j\in I}k_j^{n_j/n_I},\qquad
 m_I=\prod_{j\in I}m_j^{n_j/n_I}.
 \label{fus:sample-retention}
\end{equation}
Weighted log averages have variance
$\sum_j(n_j/n_I)^2/n_j=1/n_I$. The same calculation for the joint
reference gives $r_{n_I}G_{n_I}$. Conditional log observation noise has
variance $\sum_j(n_j/n_I)^2/(an_j)=1/(an_I)$, so the evaluator itself
is inherited. Physical samples are first returned through $U_j^{-1}$ and
the retained sample is expressed in a chosen parent frame $U_I$.

\section{Messages and state-dependent scalar links}

Define a normalized message kernel by
\begin{equation}
 B_n(k,dw)=\operatorname{LN}(\log k,1-1/n)(dw),
 \label{fus:message-kernel}
\end{equation}
interpreting $B_1(k,dw)=\delta_k(dw)$. It restores unit log variance:
$q_nB_n=\operatorname{LN}(\mu^q,1)$ and
$s_nB_n=\operatorname{LN}(\mu^s,1)$. These messages give all multiplicities
a common comparison family. The kernel has a statistical meaning. For $n$
independent unit children with log samples $N(\mu,1)$, the mean log sample is
sufficient for $\mu$ and has the law of the parent's log sample; the
conditional law of one child's log sample given that mean is
$N(\bar Z,1-1/n)$, which does not depend on $\mu$. Thus $B_n$ is the
conditional law of one child sample given the parent sufficient statistic
\cite{AyJostLeSchwachhofer2015}, and this parameter independence is why fresh
messages descend exactly.

Model prediction also commutes with this interface. Put
$h_G=(c/a)^2+1/a$ and define
$B_n^{h_G}(k,dw)=\operatorname{LN}(\log k,h_G(1-1/n))(dw)$, again using the
identity at $n=1$. Both routes in the kernel identity
\begin{equation}
 G_nB_n^{h_G}=B_nG_1
 \label{fus:evaluator-square}
\end{equation}
have conditional log mean $(c/a)\log m$ and variance
$h_G-(c/a)^2/n$. Integrating against $s_n$ gives
\begin{equation}
 (s_nB_n)G_1=(s_nG_n)B_n^{h_G}
 =\operatorname{LN}\left(\frac ca\mu^s,h_G\right).
 \label{fus:predictive-square}
\end{equation}
Thus a parent can emit the original predictive message by either normalized
route.

For each sector $\alpha\in\{q,s\}$, choose a fixed $0<\chi_\alpha<1$ and set
\begin{equation}
 L_{ij}^\alpha=\chi_\alpha(\mu_i^\alpha-\mu_j^\alpha),\qquad
 \Theta_{ij}^\alpha=U_i\exp(L_{ij}^\alpha)U_j^{-1}.
 \label{fus:links}
\end{equation}
Under a local change of frame $U_i\mapsto h_iU_i$, the physical transport
changes to $h_i\Theta_{ij}^\alpha h_j^{-1}$. The reference means and
$L_{ij}^\alpha$ remain unchanged. These links are pure gauge. With
$V_i^\alpha=U_i\exp(\chi_\alpha\mu_i^\alpha)$ one has
$\Theta_{ij}^\alpha=V_i^\alpha(V_j^\alpha)^{-1}$, so every cycle product
is trivial and $\Theta_{ji}^\alpha=(\Theta_{ij}^\alpha)^{-1}$: the
transport is a law-dependent change of frame of the trivial connection.

Transport of the reference message from $j$ shifts its log mean to
$\mu_j^\alpha+\chi_\alpha(\mu_i^\alpha-\mu_j^\alpha)$. Consequently the directed
comparison cost between restored messages
\begin{equation}
 d_{ij}^\alpha=
 \frac{\sqrt{2\KL\bigl(
 \operatorname{LN}(\mu_i^\alpha,1)\|
 (\exp L_{ij}^\alpha)_{\#}\operatorname{LN}(\mu_j^\alpha,1)
 \bigr)}}{1-\chi_\alpha}
 =|\mu_i^\alpha-\mu_j^\alpha|
 \label{fus:distance}
\end{equation}
is gauge invariant. The normalized cost does not depend on $\chi_\alpha$, so
the potential and the flow below are the same for every
$\chi_\alpha\in(0,1)$, including the limit of identity transport. The
transport therefore carries no gauge content beyond passive covariance in
this example; a transport whose parameter survives the normalization would
need, for example, a noncommuting generator on a vector log space. The
equality includes the dependence of the link on both endpoint laws. For
unequal means, differentiating its log-mean residual gives
\begin{equation}
 \partial_{\mu_i^\alpha}
 \bigl[\mu_i^\alpha-\mu_j^\alpha-L_{ij}^\alpha\bigr]=1-\chi_\alpha.
\end{equation}
This factor cancels the denominator in \cref{fus:distance}. Holding the
link fixed while differentiating would replace the coupling $\lambda$
below by $\lambda/(1-\chi_\alpha)$.

\section{The continuous subgradient dynamics}
\label{sec:fus-dynamics}

For multiplicities $n_i$ with total $N=\sum_i n_i$, choose $\lambda>0$ and
the potential
\begin{equation}
 \mathcal F(\boldsymbol\mu)
 =\sum_i n_i\bigl[C+\mathcal U(\mu_i^q,\mu_i^s)\bigr]
 +\lambda\sum_{i<j}n_in_j
 \bigl(|\mu_i^q-\mu_j^q|+|\mu_i^s-\mu_j^s|\bigr).
 \label{fus:potential}
\end{equation}
All contacts are fixed and present. If source rows are recorded, take
$\beta_{ij}=\gamma_{ij}=n_j/N$, including self entries. Summing child
entries gives exactly the same parent row. Self comparison costs vanish,
while their row mass records the corresponding contact multiplicity.
There is no Gibbs row adaptation. Relative to \cref{def:dyn-potential},
the local term is $\sum_iE_{0,n_i}$ up to the constant $\sum_i(n_i-1)C$.
The peer term $\sum_j\beta_{ij}D_{ij}^\alpha$ of agent $i$ is modified in
four constitutive ways. The divergence is replaced by its normalized square
root $d_{ij}^\alpha$ of \cref{fus:distance}; the receiver's law $q_i$ is
replaced by its restored message $q_iB_{n_i}=\operatorname{LN}(\mu_i^\alpha,1)$
before the comparison; the term is multiplied by
$\tfrac12\lambda Nn_i$, so that it becomes
$\tfrac12\lambda Nn_i\sum_j\beta_{ij}d_{ij}^\alpha$; and the rows are fixed.
The restoration matters for $n_i\geq2$: with the unrestored receiver law
$\operatorname{LN}(\mu_i^\alpha,1/n_i)$ one has
$2D_{ij}^\alpha=(1/n_i-1+\log n_i)+(1-\chi_\alpha)^2(\mu_i^\alpha-\mu_j^\alpha)^2$,
whose constant is positive, so its square root would be smooth at
coincidence and would depend on $\chi_\alpha$, and parents of multiplicity at
least two could not stick. The local term has the stated KL origin; the flow
below is a declared gradient flow, not inference in a normalized model, and
no fixed sample-model ELBO identity is asserted for the entire potential.

With $W=\operatorname{diag}(n_1,n_1,n_2,n_2,\ldots)$, define the flow on
$\R^{2r}$, for $r$ current agents, by
\begin{equation}
 W\dot{\boldsymbol\mu}\in-\partial\mathcal F(\boldsymbol\mu),
 \label{fus:inclusion}
\end{equation}
where $\partial\mathcal F$ is the subdifferential of $\mathcal F$ as a
function of the whole configuration. Solutions are absolutely continuous in
the external time $t\geq0$ and satisfy this inclusion almost everywhere.
There are no boundaries or stochastic driving terms. Both law sectors have
the same basic mobility; the factor $1/n_i$ comes from their inherited Fisher
metric. In coordinates,
\begin{align}
 \dot\mu_i^q&=-a\mu_i^q+c\mu_i^s
 -\lambda\sum_{j\ne i}n_j\epsilon_{ij}^q,\\
 \dot\mu_i^s&=-a\mu_i^s+c\mu_i^q
 -\lambda\sum_{j\ne i}n_j\epsilon_{ij}^s,
 \label{fus:coordinate-flow}
\end{align}
where $\epsilon_{ij}^\alpha=-\epsilon_{ji}^\alpha$ belongs to
$\operatorname{Sign}(\mu_i^\alpha-\mu_j^\alpha)$ and
$\operatorname{Sign}(0)=[-1,1]$. Every pair uses one antisymmetric selection;
letting each agent choose its own element of $\operatorname{Sign}(0)$, as a
product of partial subdifferentials would, admits spurious solutions, for
example two coincident unit agents that move together under any constant
selection, and is not the Filippov regularization. The multiplicity cancels
from an agent's own velocity, which depends only on the multiplicities of the
others. Each agent performs overdamped planar motion, in the coordinates
$(\mu_i^q,\mu_i^s)$, in a harmonic trap under the pair potential
$n_in_j\|x_i-x_j\|_1$, so a multiplicity acts as a friction weight and as an
interaction charge; it is not an inertia.

\lemmaheading{Global well-posedness, contraction, and persistence}{fus:uniqueness}
For every multiplicity vector and every initial datum, \cref{fus:inclusion}
has a unique absolutely continuous solution on $[0,\infty)$. Two solutions
with the same multiplicities obey
\begin{equation}
 \|\boldsymbol\mu(t)-\widetilde{\boldsymbol\mu}(t)\|_W
 \leq e^{-(a-c)t}
 \|\boldsymbol\mu(0)-\widetilde{\boldsymbol\mu}(0)\|_W.
 \label{fus:contraction}
\end{equation}
Synchronized unit-multiplicity children remain synchronized in both sectors.
\begin{proof}
The potential $\mathcal F$ is convex and finite on $\R^{2r}$, so
$W^{-1}\partial\mathcal F$ is a maximal monotone operator in the $W$ inner
product. The theory of maximal monotone operators gives a unique global
absolutely continuous solution for every initial datum, whose right
derivative is the element of $-W^{-1}\partial\mathcal F(\boldsymbol\mu)$ of
minimal norm in the $W$ norm \cite{Brezis1973}; the Euclidean minimal
element would in general not stick coincident agents of unequal
multiplicities. The quadratic part is $(a-c)$-strongly convex in the $W$
metric and the absolute-value part is convex. For subgradients
$\xi,\widetilde\xi$ at two configurations, monotonicity gives
\begin{equation}
 \langle\xi-\widetilde\xi,
 \boldsymbol\mu-\widetilde{\boldsymbol\mu}\rangle
 \geq(a-c)\|\boldsymbol\mu-\widetilde{\boldsymbol\mu}\|_W^2.
\end{equation}
Substitution from \cref{fus:inclusion} and integration proves
\cref{fus:contraction}. Swapping two equal unit-multiplicity children
preserves the initial configuration, potential, and metric. Uniqueness
therefore preserves their equality. This argument concerns the full
two-sector system, not a restricted scalar equation.
\end{proof}

Because the subdifferential of a convex function is its Clarke generalized
gradient, \cref{fus:inclusion} coincides with the Filippov regularization of
the almost-everywhere gradient field $-W^{-1}\nabla\mathcal F$
\cite{Filippov1988,Clarke1983,PadenSastry1987}, so no ambiguity of solution
concept arises. Backward uniqueness fails: the flow is a contraction
semigroup and not a group, which is what permits finite-time sticking.

\propositionheading{Complete synchronization in finite time}{fus:sync}
Fix any multiplicity vector and any initial datum. Put
$R_0=\|\boldsymbol\mu(0)\|_W$, $\kappa=a-c$, and let
$D^\alpha=\max_i\mu_i^\alpha-\min_i\mu_i^\alpha$ be the spread of sector
$\alpha$. Then $\max_{i,\alpha}|\mu_i^\alpha(t)|\leq e^{-\kappa t}R_0$, and
there is a time
\begin{equation}
 T_*\leq\frac1\kappa\log\left(1+\frac{\kappa\max_\alpha D^\alpha(0)}{\lambda N}\right)
 \label{fus:sync-time}
\end{equation}
such that all agents have equal means in both sectors for every
$t\geq T_*$. Every solution therefore produces a finite sequence of
synchronization events, at most $r-1$ of them, ending in a single block.
\begin{proof}
The zero configuration is stationary, because $0\in\partial\mathcal F(0)$.
\Cref{fus:contraction} with this stationary solution gives
$\|\boldsymbol\mu(t)\|_W\leq e^{-\kappa t}R_0$, and
$|\mu_i^\alpha|\leq\|\boldsymbol\mu\|_W$ because $n_i\geq1$. Fix a sector
$\alpha$ and write $\alpha'$ for the other sector. The spread $D^\alpha$ is
absolutely continuous. At almost every $t$ all coordinates and both extrema
are differentiable. At such a time every maximizer $i$ has
$\dot\mu_i^\alpha$ equal to the derivative of the maximum, because the
maximum minus $\mu_i^\alpha$ is nonnegative and vanishes at $t$. Let $T_+$
be the set of maximizers, with multiplicity $n_+$, and average the
coordinate equation over $T_+$ with weights $n_i/n_+$. Antisymmetry of
$\epsilon_{ij}^\alpha$ cancels the internal terms, and
$\epsilon_{ij}^\alpha=1$ for $j\notin T_+$. Hence
\begin{equation}
 \frac{d}{dt}\max_i\mu_i^\alpha
 =-a\max_i\mu_i^\alpha+c\overline{\mu^{\alpha'}}_{T_+}-\lambda(N-n_+),
\end{equation}
where $\overline{\mu^{\alpha'}}_{T_+}$ is the weighted average over $T_+$.
The same argument for the minimizing set $T_-$, with multiplicity $n_-$,
gives the opposite sign of the last term. While $D^\alpha>0$ the sets $T_+$
and $T_-$ are disjoint, so $n_++n_-\leq N$, and both weighted averages of
$\mu^{\alpha'}$ lie in the range of the $\alpha'$ means. Therefore
\begin{equation}
 \dot D^\alpha\leq-aD^\alpha+cD^{\alpha'}-\lambda N
 \quad\text{almost everywhere on }\{D^\alpha>0\}.
\end{equation}
Let $M=\max(D^q,D^s)$. At almost every $t$ where $M>0$, the derivative of
$M$ equals $\dot D^\alpha$ for a sector $\alpha$ with $D^\alpha=M$ (where
$D^q=D^s$ the two derivatives agree almost everywhere), so
$\dot M\leq-\kappa M-\lambda N$ almost everywhere on $\{M>0\}$. Hence
$e^{\kappa t}M(t)\leq M(0)-\lambda N(e^{\kappa t}-1)/\kappa$ as long as $M$
stays positive, so $M$ vanishes by the time \cref{fus:sync-time}. Agents
that coincide in both sectors remain coincident: expanding integer
multiplicities into unit children and applying \cref{fus:uniqueness} proves
this, so $M$ stays zero afterward. The number of distinct agents therefore
decreases at each event.
\end{proof}

The bound depends only on the initial spreads; it does not involve $R_0$ or
$c$, and it is attained by two blocks on the common-offset family below. The
common-offset family studied next provides explicit event times; it is not
needed for finite-time synchronization. The following estimate controls a
single pair off that family.

\lemmaheading{Isolated pairs synchronize}{fus:isolated-pair}
Let $i\neq j$ and let $\mathcal T$ be a time interval on which, in each
sector $\alpha$, no third agent has a mean in the closed interval between
$\mu_i^\alpha$ and $\mu_j^\alpha$. Put
$S=|\mu_i^q-\mu_j^q|+|\mu_i^s-\mu_j^s|$. Then, for $s<t$ in $\mathcal T$,
\begin{equation}
 S(t)\leq\max\{0,S(s)-\lambda(n_i+n_j)(t-s)\}.
 \label{fus:pair-decay}
\end{equation}
\begin{proof}
For a third agent $k$, the hypothesis gives
$\epsilon_{ik}^\alpha=\epsilon_{jk}^\alpha$, so its terms cancel from the
equation for $\Delta^\alpha=\mu_i^\alpha-\mu_j^\alpha$, leaving
$\dot\Delta^\alpha=-a\Delta^\alpha+c\Delta^{\alpha'}
-\lambda(n_i+n_j)\epsilon_{ij}^\alpha$. Where $\Delta^\alpha\neq0$,
$\frac{d}{dt}|\Delta^\alpha|\leq-a|\Delta^\alpha|+c|\Delta^{\alpha'}|
-\lambda(n_i+n_j)$; where $\Delta^\alpha=0$ on a set of positive measure,
the derivative of $|\Delta^\alpha|$ vanishes almost everywhere. Adding the
two sectors gives $\dot S\leq-\kappa S-2\lambda(n_i+n_j)$ when both
differences are nonzero, and $\dot S\leq-\lambda(n_i+n_j)$ when exactly one
is nonzero. Hence $\dot S\leq-\lambda(n_i+n_j)$ almost everywhere on
$\{S>0\}$, which proves \cref{fus:pair-decay}.
\end{proof}

\section{An invariant class with explicit global solutions}

\theoremheading{Finite sticking on a common-offset family}{fus:sticking}
Let all initial multiplicities be one and suppose
\begin{equation}
 \mu_i^q(0)=u_i^0+v_0,\qquad
 \mu_i^s(0)=u_i^0-v_0,
 \label{fus:admitted-class}
\end{equation}
where the $u_i^0$ are distinct real numbers and $v_0\ne0$. The unique
global solution of the full inclusion for these initial data has, in both
sectors, the form
\begin{equation}
 \mu_i^q=u_i+v,\qquad \mu_i^s=u_i-v,\qquad
 v(t)=v_0e^{-(a+c)t}.
 \label{fus:invariant-form}
\end{equation}
Blocks formed by equal $u_i$ persist and follow the multiplicity-weighted
version of the same equation.
\begin{proof}
Set $\kappa=a-c>0$. Equal sector orderings, with the same antisymmetric
selection $\epsilon_{ij}=\epsilon_{ij}^q=\epsilon_{ij}^s$ in both sectors,
reduce \cref{fus:coordinate-flow} to
\begin{equation}
 \dot u_i=-\kappa u_i-\lambda\sum_{j\ne i}n_j\epsilon_{ij},\qquad
 \epsilon_{ij}=-\epsilon_{ji}\in\operatorname{Sign}(u_i-u_j),\qquad
 \dot v=-(a+c)v.
\end{equation}
Introduce the increasing auxiliary time and rescaled positions
\begin{equation}
 \tau(t)=\frac{e^{\kappa t}-1}{\kappa},\qquad
 u_i(t)=e^{-\kappa t}z_i(\tau(t)).
 \label{fus:time-change}
\end{equation}
Since $\dot\tau=e^{\kappa t}>0$, direct differentiation gives
\begin{equation}
 z_i'=-\lambda\sum_{j\ne i}n_j\epsilon_{ij},\qquad
 \epsilon_{ij}=-\epsilon_{ji}\in\operatorname{Sign}(z_i-z_j).
 \label{fus:sticky-flow}
\end{equation}
This is the sticky-particle solution of the one-dimensional aggregation
equation with kernel $\lambda|z|$ on weighted atoms
\cite{CarrilloEtAl2011,BonaschiEtAl2015}. An ordered block has constant
velocity $\lambda(n_{\rm right}-n_{\rm left})$ between collisions, where
$n_{\rm right}$ and $n_{\rm left}$ are the total multiplicities to its
right and left. Adjacent blocks approach with relative speed
$\lambda(n_I+n_J)>0$. At their collision, join their multiplicities, retain
the common position, and use the new velocity
$\lambda(n_{\rm right}-n_{\rm left})$. This equals the multiplicity-weighted
average of their velocities just before collision. Choosing zero internal
subgradients verifies the original inclusion for every constituent child
after collision. Each event reduces the number of blocks, so at most
$N-1$ events occur, and the resulting piecewise-linear curve exists for
all $\tau\geq0$. Reversing \cref{fus:time-change} constructs a global
full-system solution. \Cref{fus:uniqueness} makes it the solution.
\end{proof}

The class \cref{fus:admitted-class} has a common difference between its
two sector means. This is an initial-data restriction, not an imposed
constraint force. It is an invariant affine subspace of codimension $r-1$
in $\R^{2r}$. The two laws of each agent remain distinct at every finite
time because $v(t)\ne0$. The passage to $(z,\tau)$ is a rescaling of
positions by the common positive factor $e^{\kappa t}$ combined with a
monotone change of time; it preserves coincidence, and $\tau$ is not a
geometric or operational clock.

The merge order on this family has a closed form. For a set $G$ of
consecutive sorted agents, write $n_G$ for its multiplicity and $c_G$ for
the multiplicity-weighted mean of the initial offsets $u_i^0$, $i\in G$.
The normalized-centroid linkage of the initial offsets starts from the
singletons and repeatedly merges the pair of adjacent current blocks
$G<H$ that minimizes
\begin{equation}
 \varrho_{GH}=\frac{c_H-c_G}{n_G+n_H}.
 \label{fus:ncl-height}
\end{equation}

\propositionheading{The coalescence tree is a coupling-independent linkage of the initial data}{fus:ncl}
In \cref{fus:sticking}, allowing arbitrary initial multiplicities, the
membership history consists exactly of the merges of the normalized-centroid
linkage of the initial offsets, in its order, and the $k$-th merge, of
height $\varrho_k$, occurs at
\begin{equation}
 \tau_k=\frac{\varrho_k}\lambda,\qquad
 t_k=\frac1\kappa\log\left(1+\frac{\kappa\varrho_k}{\lambda}\right).
 \label{fus:ncl-times}
\end{equation}
The heights are nondecreasing. Consequently the partition sequence is the
same for every $\lambda>0$ and every $a>c>0$; only the event times depend on
these constants.
\begin{proof}
Work in the $z$ frame of \cref{fus:sticky-flow}, where the order of the
agents is preserved and contact is sticky. Let $G$ be a current block at
auxiliary time $\tau$. No agent outside $G$ has touched a member of $G$
before $\tau$, since by stickiness it would then belong to $G$. Hence, at
every earlier time, the signs between members of $G$ and outsiders are single
valued and constant, the internal terms cancel by antisymmetry, and
$\frac{d}{d\tau}\sum_{i\in G}n_iz_i=\lambda n_G(n_{\rm right}(G)-n_{\rm left}(G))$
from $\tau=0$. For adjacent current blocks $G<H$ one has
$n_{\rm left}(H)=n_{\rm left}(G)+n_G$ and $n_{\rm right}(G)=n_{\rm right}(H)+n_H$,
so the distance of their weighted centroids, which is their gap, equals
$(c_H-c_G)-\lambda(n_G+n_H)\tau$. They therefore meet at
$\tau=\varrho_{GH}/\lambda$ unless an earlier event involves one of them, and
the next event is a merge of the adjacent current pair with the smallest
$\varrho_{GH}$. Induction on the number of blocks proves the first claim.
When $G$ and $H$ merge into $K$ at $\tau_k$, the gap between $K$ and a
neighbor $L$ is still nonnegative at $\tau_k$, so
$\varrho_{KL}\geq\lambda\tau_k=\varrho_k$, which gives monotonicity.
Inverting \cref{fus:time-change} gives $t_k$. The heights are computed from
the initial offsets and multiplicities alone.
\end{proof}

\section{Two fusion stages without a supplied partition}

\theoremheading{Four agents fuse into two parents and then one}{fus:two-stages}
In \cref{fus:sticking}, take four agents with the unordered initial set
\begin{equation}
 \{u_i^0\}=\{-3d,-2d,2d,3d\},\qquad d>0.
\end{equation}
Two size-two blocks form at $t_1$ and remain distinct until $t_2>t_1$, where
they fuse under their inherited equations, with
\begin{equation}
 t_1=\frac1\kappa\log\left(1+\frac{\kappa d}{2\lambda}\right),\qquad
 t_2=\frac1\kappa\log\left(1+\frac{5\kappa d}{4\lambda}\right).
 \label{fus:collision-times}
\end{equation}
\begin{proof}
The adjacent heights \cref{fus:ncl-height} of the singletons are $d/2$,
$2d$, and $d/2$, so the two outer pairs merge first, at height $d/2$. The
two resulting blocks have centroids $\mp5d/2$ and multiplicity two, so their
height is $5d/4$. \Cref{fus:ncl} gives \cref{fus:collision-times}.
Equivalently, in sorted order the initial $z$ velocities are
$(3\lambda,\lambda,-\lambda,-3\lambda)$, the outer gaps close at
$\tau_1=d/(2\lambda)$ at the parent positions $\mp3d/2$, and the parents,
with velocities $\pm2\lambda$, meet after another $3d/(4\lambda)$.
\end{proof}

Exact reflection symmetry is unnecessary. For four sorted initial positions,
let their adjacent gaps be $g_1,g_2,g_3>0$, and put $H=\max(g_1,g_3)$. The
two outer pairs form first and then fuse if and only if
\begin{equation}
 \Gamma=g_2+\frac{g_1+g_3}{2}-2H>0.
 \label{fus:open-gap-condition}
\end{equation}
Suppose $g_1\leq g_3$; the other case is the reflected calculation. Before
the first outer collision the middle gap decreases at speed $2\lambda$, so
at $g_1/(2\lambda)$ it equals $g_2-g_1$; afterward it decreases at speed
$3\lambda$ until the other outer pair collides at $g_3/(2\lambda)$, where
the same linear expression has the value $\Gamma$. If $\Gamma>0$, then
$g_2>H$, the middle gap stays positive until both outer pairs have formed by
$\tau_*=H/(2\lambda)$, and the two parent gap then decreases at speed
$4\lambda$, so the final time is
\begin{equation}
 \tau_{\rm final}=\tau_*+\frac{\Gamma}{4\lambda}
 =\frac{2g_2+g_1+g_3}{8\lambda}.
\end{equation}
If $\Gamma\leq0$, the middle gap closes no later than the second outer
collision, and the pattern fails; at $\Gamma=0$ all four agents merge in one
event, as three current blocks if $g_1\neq g_3$ and as four singletons if
$g_1=g_3$. In the language of \cref{fus:ncl}, $\Gamma>0$ is the statement
that the normalized-centroid linkage merges the two outer pairs before
either joins the middle. Condition \cref{fus:open-gap-condition} is strict
and defines an open set relative to the ordered common-offset family. Taking
all label permutations gives an admitted class with no assigned labels for
blocks. The initial laws are unequal across agents, and the interaction rule
contains no partition. Membership is selected by attained equality.

\corollaryheading{The two-stage pattern is open in the full space}{fus:open-two-stage}
Let $\bar x\in\R^8$ be a four-agent configuration of unit multiplicities in
the common-offset class satisfying \cref{fus:open-gap-condition}, and let
$\mathcal O_1$ and $\mathcal O_2$ be its two outer pairs. There is
$\varepsilon>0$ such that every initial configuration within distance
$\varepsilon$ of $\bar x$ has the following membership history:
$\mathcal O_1$ and $\mathcal O_2$ synchronize in both sectors, in some order
and before any other synchronization, and the two resulting parents then
synchronize. The set of initial data with this history therefore contains an
open subset of $\R^8$ that contains the open-gap set. Explicitly, with the
gaps ordered so that $g_1\leq g_3$ and $t_{\rm f}$ the final time of
$\bar x$, every $\varepsilon<\Gamma e^{-\kappa t_{\rm f}}/16$ suffices, and
no radius larger than $(2\sqrt2/5)\Gamma$ does.
\begin{proof}
Let $\bar{\boldsymbol\mu}$ be the solution from $\bar x$, let $t_*$ be its
second outer collision time and $t_{\rm f}>t_*$ its final time, and choose
$\eta>0$ with $t_*+\eta<t_{\rm f}$. On $[0,t_*+\eta]$ every mean of
$\mathcal O_1$ lies below every mean of $\mathcal O_2$ in each sector, with a
margin $\sigma_0>0$. By \cref{fus:contraction}, a solution from a
configuration within $\varepsilon$ of $\bar x$ stays within $\varepsilon$
of $\bar{\boldsymbol\mu}$ in every coordinate. If $2\varepsilon<\sigma_0$,
both pairs are isolated in the sense of \cref{fus:isolated-pair} on
$[0,t_*+\eta]$, and no agent of $\mathcal O_1$ meets an agent of
$\mathcal O_2$ in either sector there. At the collision time of
$\mathcal O_1$ along $\bar{\boldsymbol\mu}$, the perturbed pair has
$S\leq4\varepsilon$, so \cref{fus:pair-decay} synchronizes it within
$2\varepsilon/\lambda$; the same holds for $\mathcal O_2$. If also
$2\varepsilon/\lambda\leq\eta$, both pairs synchronize before
$t_*+\eta$ and no other event occurs before then. \Cref{fus:uniqueness}
keeps each pair synchronized, and \cref{fus:sync} synchronizes the two
parents at a finite later time.

For the explicit radius, choose $\eta$ by
$\tau(t_*+\eta)=\tau_*+\Gamma/(8\lambda)$, which is below $\tau_{\rm final}$.
Along $\bar{\boldsymbol\mu}$ the gap between the two pairs is
$e^{-\kappa t}$ times the $z$ gap, a product of two decreasing positive
factors, and at $\tau_*+\Gamma/(8\lambda)$ the $z$ gap is $\Gamma/2$; hence
$\sigma_0\geq e^{-\kappa t_{\rm f}}\Gamma/2$. Since
$\eta=\kappa^{-1}\log\bigl(1+\kappa\Gamma/(8\lambda(1+\kappa\tau_*))\bigr)$
and $\log(1+y)\geq y/(1+y)$, one has
$\eta\geq\Gamma/(8\lambda(1+\kappa\tau_{\rm final}))=\Gamma e^{-\kappa t_{\rm f}}/(8\lambda)$.
Both requirements hold when $\varepsilon<\Gamma e^{-\kappa t_{\rm f}}/16$.
Conversely, lowering the third sorted offset by $\delta$ in both sectors
changes $(g_2,g_3)$ to $(g_2-\delta,g_3+\delta)$ and $\Gamma$ to
$\Gamma-5\delta/2$, which vanishes at $\delta=2\Gamma/5$; that configuration
lies at distance $(2\sqrt2/5)\Gamma$ from $\bar x$ and has a four-agent
merger.
\end{proof}

The admissible radius is therefore of order $\Gamma$, and the open set
shrinks toward the boundary $\{\Gamma=0\}$ of the open-gap set.

\corollaryheading{Coalescence of any finite depth}{fus:dyadic}
Place $N=2^k$ unit agents in the common-offset class with sorted initial
positions $u^0$, and assign to the gap between sorted agents $p$ and $p+1$
the level of the smallest dyadic block of indices containing both. If, for
each level $\ell<k$,
\begin{equation}
 2^\ell\max\{\text{gaps of level }\ell\}
 <\min\{\text{gaps of level greater than }\ell\},
 \label{fus:dyadic-condition}
\end{equation}
then all dyadic blocks of size $2^\ell$ form before any event that joins two
of them, for $\ell=1,\ldots,k$, and each stage is inherited exactly. The
condition is strict, so the $k$-stage history holds on an open subset of the
common-offset class.
\begin{proof}
In the $z$ frame, adjacent blocks with multiplicities $n_I,n_J$ approach at
speed $\lambda(n_I+n_J)\geq2\lambda$, so a gap of initial size $g$ between
two sorted agents closes by the auxiliary time $g/(2\lambda)$. By
\cref{fus:dyadic-condition} at the lower levels, the maximal gap of level
$\ell$ exceeds every gap of lower level, so every gap of level at most
$\ell$ has closed by $\max\{\text{gaps of level }\ell\}/(2\lambda)$. Before
the first closure of a gap of level greater than $\ell$, every current block
lies inside a dyadic block of size $2^\ell$, so adjacent blocks approach at
speed at most $\lambda2^{\ell+1}$, and no gap of level greater than $\ell$
closes before $\min\{\text{gaps of level greater than }\ell\}/(\lambda2^{\ell+1})$,
which by \cref{fus:dyadic-condition} is later. Induction on $\ell$ proves
the ordering. Event times are $t=\kappa^{-1}\log(1+\kappa\tau)$, and
inheritance at every stage is the handoff identity of the next section.
\end{proof}

The condition $\frac N2\max\{\text{gaps of level }\ell\}<\min\{\text{gaps of
level }\ell+1\}$ for every $\ell$ implies \cref{fus:dyadic-condition},
because $N/2\geq2^\ell$ and the inequalities chain across levels. Every such
hierarchy ends in one block, by \cref{fus:sync}.

\section{The inherited handoff and its scope}

At a collision, restrict a block $I$ to its common means. Summing its local
energies and pulling back its metric gives
\begin{equation}
 \sum_{i\in I}n_i(C+\mathcal U)
 =n_I(C+\mathcal U)=E_{0,n_I}+(n_I-1)C,\qquad
 g_I=n_I\bigl((d\mu_I^q)^2+(d\mu_I^s)^2\bigr).
\end{equation}
Every external pair term sums to
$n_In_J|\mu_I^\alpha-\mu_J^\alpha|$. Zero internal subgradients therefore
lift each solution of the inherited equation to a fine solution. Uniqueness
identifies that lift with the original trajectory. This also proves
persistence for inherited multiplicities by expanding a multiplicity into
synchronized unit children. With $R$ retaining block means, multiplicities,
and frames, the corresponding parameter flows satisfy the following
identity. For a fixed formed partition, keep both component index sets
fixed even when further components coincide; a later merger is a separate
retention composed with $R$. Then
\begin{equation}
 R\circ\Phi_t=\Phi_t^c\circ R
 \label{fus:flow-square}
\end{equation}
on every synchronized stratum reached by the construction, for all future
times. Weighted geometric averaging is associative, so later collisions
give the same retained laws as direct averaging. The state is continuous
at a collision; the velocity changes through the inherited multiplicity sum.

The averaged dynamics alone does not single out synchronized blocks. Call a
partition separated if, in each sector and for every pair of blocks, all
means of one block lie strictly below all means of the other. Then every
external sign is single valued and constant, the internal terms cancel by
antisymmetry, and the weighted block averages obey the coarse equation until
two blocks interleave, whether or not the blocks are synchronized. What
holds only on synchronized strata is forward invariance of the retained
structure and inheritance of the declared responses: for unsynchronized
children the log samples given the retained log mean have a law that depends
on the discarded deviations, so fresh child messages and the joint
reconstruction below do not descend. The selected partition is therefore
distinguished by response inheritance and invariance, not by the averaged
dynamics.

Along every solution the convex chain rule for absolutely continuous curves
\cite{Brezis1973} gives
\begin{equation}
 \frac{d\mathcal F}{dt}
 =-\sum_i n_i\bigl((\dot\mu_i^q)^2+(\dot\mu_i^s)^2\bigr)
 \quad\text{almost everywhere},
\end{equation}
including intervals on which pairs slide together in one sector only. The
objective is continuous at collisions when the constants above are
retained. By \cref{fus:sync}, every solution converges to the stationary
configuration with every mean zero, after complete synchronization at a
finite time.

\remarkheading{A discrete update that keeps exact sticking}{fus:prox}
The natural time discretization with step $\Delta t>0$ is the proximal step
\begin{equation}
 \boldsymbol\mu^{k+1}=\operatorname*{argmin}_{x}
 \left\{\mathcal F(x)+\frac1{2\Delta t}\|x-\boldsymbol\mu^k\|_W^2\right\},
 \label{fus:prox-step}
\end{equation}
whose minimizer is unique because the objective is strongly convex. Its
optimality condition
$W(\boldsymbol\mu^k-\boldsymbol\mu^{k+1})/\Delta t\in\partial\mathcal F(\boldsymbol\mu^{k+1})$
and strong monotonicity give the discrete contraction
$\|\boldsymbol\mu^{k+1}-\widetilde{\boldsymbol\mu}^{k+1}\|_W\leq
(1+\kappa\Delta t)^{-1}\|\boldsymbol\mu^k-\widetilde{\boldsymbol\mu}^k\|_W$,
and comparing the objective at its minimizer with its value at
$\boldsymbol\mu^k$ gives
$\mathcal F(\boldsymbol\mu^{k+1})+(2\Delta t)^{-1}\|\boldsymbol\mu^{k+1}-\boldsymbol\mu^k\|_W^2
\leq\mathcal F(\boldsymbol\mu^k)$. Uniqueness and the exchange symmetry of
equal-multiplicity agents keep synchronized blocks synchronized, and the
step creates exact ties: for two unit agents under the pair term alone, the
step returns equal means exactly when the gap before the step is at most
$2\lambda\Delta t$. The iterates converge to the solution of
\cref{fus:inclusion} uniformly on bounded time intervals as $\Delta t\to0$
\cite{CrandallLiggett1971}. An explicit Euler step with sign selections
keeps ties that are already present but does not create exact ties
generically; it oscillates around coincidence, so equality-based membership
is only meaningful for the proximal scheme or for an event-driven
integration of \cref{fus:sticky-flow}.

Sample retention is genuinely noninjective: different positive sample
vectors have the same geometric mean. For $n$ synchronized unit children,
their log samples conditional on the retained log mean $\bar Z$ have law
\begin{equation}
 (Z_1,\ldots,Z_n)\mid\bar Z
 \sim N\left(\bar Z\mathbf1,
 I_n-\frac{\mathbf1\mathbf1^{\mathsf T}}n\right).
 \label{fus:reconstruction}
\end{equation}
The residual is independent of the mean: their joint Gaussian cross
covariance is zero. Exponentiation gives a normalized joint reconstruction
kernel. Its one-child marginal is \cref{fus:message-kernel}. Independent
residuals for all children would be incorrect because the off-diagonal
conditional covariances are $-1/n$. Generative joint samples similarly use
the covariance of $(\log m,\log k)$ under $s_1G_1$; their
residual kernel is determined by that covariance, not by the recognition
product law. The local KL residual $(n-1)C$ records the corresponding
state-independent information loss under the reference comparison.

\remarkheading{Sufficiency beyond the lognormal family}{fus:sufficiency}
The kernel \cref{fus:reconstruction} is the conditional law of the
children given their sufficient statistic, and it is free of the parameter
for that reason. For an exponential family with sufficient statistic $T$,
the same holds for independent children of equal parameter: the parent law
is the law of $\sum_jT(X_j)$, the Fisher metric of the children restricted to
equal parameters is multiplied by the multiplicity, and the parameter-free
conditional law supplies both the message kernel and the joint
reconstruction. Under fixed support and the usual regularity, a sufficient
statistic of fixed dimension for every number of children exists only for
exponential families, by the Pitman--Koopman--Darmois theorem, so this is the
scope of the remark. For Gamma children of shape $k$, for example, one child
given the sum $S$ of $n$ children is $S$ times a
$\operatorname{Beta}(k,(n-1)k)$ variable. Two ingredients are specific to
Gaussian location families with fixed covariances. First, by the chain rule
for relative entropy, $\sum_{i\in I}E_{0,n_i}-E_{0,n_I}$ equals the expected
KL divergence between the parameter-free recognition conditional of the
children given $T$ and the conditional of the reference given $T$; for
Gaussian location families both conditionals are Gaussian with the same
mean and state-independent covariances, which gives the constant $(n_I-1)C$
above, and the residual vanishes when the joint reference belongs to an
exponential family for which $T$ is sufficient. In general it depends on the
state, for example for Bernoulli children with an evaluator proportional to
$e^{Jkm}$, and the restricted energy is then the parent's own KL plus a
state-dependent term, so the parent is not an agent of the children's type
in the sense of the local energy. Second, the square-root divergence between
restored messages is the convex function $|\Delta\mu|$; on a general regular
family it equals the Fisher distance only to leading order near the
diagonal. Global uniqueness would then need semiconvexity of the square-root
divergence along Fisher geodesics on a forward-invariant compact set; this
fails near the boundary of the Bernoulli family, its constant grows with
$\lambda$, and semiconvexity gives uniqueness but not the contraction used in
\cref{fus:sync}.

Parameter retention has a narrower meaning. On an exactly synchronized
stratum, child means are redundant copies; retaining their common value
does not identify different nonsynchronized configurations. Away from that
stratum, the weighted average closes the dynamics only on separated
partitions, as explained above. \Cref{fus:flow-square} makes no claim about
such hidden parameter fibers.

The admitted boundary protocol draws fresh messages through $B_n$, fresh
predictive messages through \cref{fus:predictive-square}, or complete
fresh represented samples through the joint reconstruction. To type an
accessible stored readout, augment the complete population configuration by
$o\in\mathsf O=(0,\infty)$ and set $\rho(\boldsymbol\mu,\boldsymbol n,o)=o$.
At declared epochs, the transition $\Gamma_{\rm rd}$ advances the
deterministic parameter flow and overwrites $o$ using the selected normalized
scalar output kernel. Between epochs, $o$ is fixed, and it never enters the
mean equation. Retention copies this accessible $o$ exactly; it does not
claim access to cached private residual samples. For vector readouts the same
definition uses the corresponding finite product space and joint kernel.
Equal parameter trajectories and conditional output kernels imply equal
joint accessible readout laws for every finite sequence of such queries,
by multiplying successive conditional kernels. Reusing persistent latent
samples or introducing record feedback requires another state specification.

\section{Status against the formation standard}
\label{sec:fus-status}

We report the construction by the scheme following
\cref{inh:recursive-formation}. The objective is the nonsmooth modification
\cref{fus:potential} of the reciprocal law potential; the dynamics is the
subdifferential inclusion \cref{fus:inclusion}, a declared restricted-family
flow on the parameter space $\R^{2r}$ of log means, well posed by
\cref{fus:uniqueness}. (F1) holds in the full grade: the admitted class is
the open subset of $\R^8$ given by \cref{fus:open-two-stage}, with all label
permutations, and its trace on the common-offset family, an invariant family
of codimension three, is open there. The selector is the partition of the
current state into agents with equal means in both sectors. It has no real
constant, so the declared family is a single selector; it reads laws only,
it is a Borel function of the state, and its persistence criterion is the
invariant block relation of \cref{fus:uniqueness}; (F2) holds. The parent
protocol rule $A_I$ consists of fresh messages through $B_{n_I}$, fresh
predictive messages through \cref{fus:predictive-square}, complete fresh
represented samples through the joint reconstruction, and the stored readout,
each addressing the block as a whole. The child-level probe rule $B_I$
consists of these protocols applied to individual members and the peer
channels applied to each member, with response laws the restored messages
of that member. On synchronized strata the laws, evaluator, metric, rows,
contacts, and every protocol of $A_I$ are inherited exactly, and the parent
flow is the multiplicity-weighted inclusion, well posed by
\cref{fus:uniqueness}, so (F3) holds; the geometric clause is not declared,
since the links are pure gauge.

Clause (F4) fails. The reachable post-handoff set consists of configurations
whose block members are identical copies, so the retention of common means
and multiplicity is injective there, and no probe of $B_I$ has a pair of
distinct configurations to separate. Child frames are fixed passive data.
The geometric mean is noninjective on sample vectors, but samples are not
configurations, and the conditional law of the child samples given the
retained sample is the parameter-free kernel \cref{fus:reconstruction}.

For the controls of \cref{inh:null-flow}, the peer costs are the pair terms
of \cref{fus:potential}, whose compared laws depend on another agent's law,
and the self costs vanish identically. Deleting the peer costs, with or
without their labels, gives the same flow, the case $\lambda=0$, so the
label-preserving and renormalized nulls coincide. Each agent then follows the
linear equation $\dot\mu_i=-A\mu_i$, which is globally well posed and
invertible, so distinct agents never coincide and $P^0$ contains no event.
By \cref{fus:sync}, every configuration with two distinct agents has at
least one event along $\Phi$, so $P\neq P^0$ at every point of the admitted
class. For the declared selector (F5) therefore holds against both nulls,
in the strong grade. This verdict measures the exactness of the selector
rather than which blocks the interaction forms. No backward-unique null flow
produces an attained equality, so every finite-time sticking interaction
read by an exact-equality selector passes in this way. On the common-offset
family the partition sequence along $\Phi$ is the normalized-centroid
linkage of the initial offsets and is the same for every $\lambda>0$
(\cref{fus:ncl}). Along both nulls every difference of means is
$e^{-\kappa t}$ times its initial value in both sectors, so a selector that
joins two adjacent current groups when the normalized centroid distance
\cref{fus:ncl-height} of their current offsets reaches $\delta$ from above
gives, for every $\delta$ below the smallest linkage height, the same
linkage sequence along the nulls; along $\Phi$ it gives that sequence as
well, since by the proof of \cref{fus:ncl} the threshold is reached for the
pair of height $\varrho$ at $t=\kappa^{-1}\log\bigl((\kappa\varrho+\lambda)/(\kappa\delta+\lambda)\bigr)$,
an increasing function of $\varrho$, before the corresponding coincidence.
For such a threshold family, with constants in an interval below the
smallest linkage height, $P=P^0$ on the family-grade admitted class, and
(F5) fails there. The interaction determines the event times and the
exactness of coincidence, not which blocks form. The selector-free
discrepancy between $\Phi$ and the nulls is the finite-time attainment of
equality; both converge to the zero configuration. The frozen-environment
control decouples the agents into single-agent subgradient flows of strongly
convex functions, well posed by the argument of \cref{fus:uniqueness}; its
partition sequence is not determined here. Whether the partition sequence
depends on $\lambda$ off the common-offset family is open.

Clause (F6a) holds: by \cref{fus:two-stages} and \cref{fus:open-two-stage},
the two first-stage parents select their union under the inherited
interaction, with the same selector and constitutive data, no reset of their
data, and no retuned coupling. Clause (F6b) fails for every counterfactual
class of \cref{inh:counterfactual-class}: by \cref{fus:sync} applied to
multiplicities two, every configuration of two first-stage parents, whatever
their means, fuses at a finite time, so the second-stage event occurs from
every counterfactual handoff, and the stages form a supplied-scale sequence
ordered by the initial gaps of \cref{fus:open-gap-condition}. The
amplifying grade holds only in the degenerate sense of
\cref{inh:recursive-formation}, since within-block diameters vanish at every
event of $\Phi$.

The construction satisfies (F1), (F2), (F3), and (F6a), and (F4) and
(F6b) are shown to fail, so it is a two-stage selection with exact
inheritance and not a recursive formation. Its (F5) verdict holds only
through the exact-equality selector; its coalescence tree is a
coupling-independent linkage of the initial data, and the tree always ends in
one block.

\remarkheading{Smooth peer KL dynamics and persistent parents}{fus:smooth-open}
The finite mergers above depend on the nonsmooth interaction. Under the
smooth reciprocal potential with a locally Lipschitz vector field, an
interior trajectory cannot first reach an invariant synchronization
manifold at finite time (\cref{inh:smooth-no-finite-entry}), so smooth
constructions need a different selection criterion; the threshold
selectors of \cref{ch:conditional-formation,ch:open-formation} are such
criteria, and the remaining problem is \cref{inh:formation-open}. Persistent
parents are excluded here by \cref{fus:sync}. A contact range, for example
$\lambda_{ij}=\lambda\mathbf 1\{|\mu_i-\mu_j|<\varrho_{\rm c}\}$ as in
bounded-confidence dynamics \cite{HegselmannKrause2002}, or a multiwell local
reference would let separated blocks stop interacting. Both destroy the
convexity on which uniqueness rests, and well-posedness and inheritance for
such variants have not been established. Persistence does not by itself
require giving up convexity: agent-specific single-well references
$n_i\mathcal U(\mu_i-x_i)$ keep the potential $(a-c)$-strongly convex, so
\cref{fus:uniqueness} and \cref{fus:contraction} hold verbatim, and the
unique minimizer can have several blocks. Under any strongly convex
potential, however, every solution converges to that minimizer whatever the
initial datum, so the limiting partition is a function of the declared data
alone and a final-stage counterfactual dependence (F6b) needs several minima.
For such references, finite-time arrival at the limiting partition, and the
behavior of blocks before it, are not established here; synchronized blocks
need not persist before arrival, since agent-specific references break the
exchange symmetry used in \cref{fus:uniqueness}.
